% ============================================================
% paper/appendices/appx_formal_skeleton.tex
% Appendix D: Formal Skeleton (Invariants / STOP / No-Go)
% Purpose: theorem-style restatement without proofs.
% ============================================================

\section*{Appendix D: Formal Skeleton (Invariants, STOP, No-Go)}

\subsection*{Purpose of This Appendix}

This appendix addresses the ``formalism without demonstration'' critique
by providing a compact theorem-style skeleton of the core structural
claims already present in the main text and in \texttt{ETS.K\_EA.v1.1}.

No proofs are provided. No new theorems are introduced.
The appendix is a restatement of invariants, precedence, and
termination outcomes as a checklist for reviewers.

\subsection*{Notation (Local to This Appendix)}

All symbols are used in the sense of the main document:
$\OmegaEA$, $\dOmegaEA$, $\AEA$, $\ThetaEA$, $\JEA$, $\CEA$, $\kEA$.

\subsection*{Skeleton Statements (Restatement Only)}

\begin{theorem}[Executable closure precedence]
If a textual statement in this whitepaper appears to contradict
\texttt{ETS.K\_EA.v1.1}, the ETS takes precedence. The whitepaper has no
authority to override executable constraints.
\end{theorem}

\begin{theorem}[STOP as definitional outcome]
STOP is an admissible outcome whenever structural existence,
observability, or compatibility conditions required by the ETS are not
met. STOP indicates that further diagnosis is undefined within the model,
not that the model has failed.
\end{theorem}

\begin{theorem}[Admissibility non-negotiability]
If $\OmegaEA=\varnothing$, the continuum $K_{\mathrm{EA}}$ has no
admissible states under the declared constraints, and diagnosis is
undefined (STOP).
\end{theorem}

\begin{theorem}[Boundary sensitivity]
Approaching $\dOmegaEA$ corresponds to approaching threshold equality or
admissibility limits. Boundary proximity does not license prescriptions;
it only delimits the defined domain of the formal object.
\end{theorem}

\begin{theorem}[Continuity measure status]
$\kEA$ is a structural measure summarizing coherence of admissibility
under declared axes, flows, cycles, and thresholds. $\kEA$ is not an
empirical KPI and is not presented as cross-enterprise comparable.
\end{theorem}

\begin{theorem}[No-Go: bypassing termination]
No figure, instantiation, or external artifact may be used to bypass
STOP conditions, invariants, compatibility constraints, or phase
precedence rules defined by the ETS.
\end{theorem}

\begin{theorem}[No-Go: importing external authority]
External literature, tools, or reconstruction artifacts have no
normative authority over the ETS. They may support inspection of the
artifact identity but cannot modify admissibility or phase logic.
\end{theorem}

\subsection*{Interpretation Boundary}

These skeleton statements are not claims of empirical truth.
They are structural consequences of a closed specification.

%%%End of Document%%%
